\documentclass[10pt,a4paper]{amsart}
\usepackage[margin=19mm]{geometry}
\usepackage{amsmath,amssymb,mathtools}
\usepackage[scr=rsfs,cal=euler]{mathalfa}
\usepackage{xfrac}
\usepackage[unicode,hidelinks,bookmarks=false]{hyperref}
\newcommand{\ie}{i.e.\ }
\newcommand{\cf}{cf.\ }
\newcommand{\resp}{respectively\ }
\newcommand{\Subsection}{\subsection}
\providecommand{\nobreakdash}{\nobreak\hbox{-}}
\setlength{\emergencystretch}{2em}
\multlinegap0pt
\usepackage{bm}
\usepackage{bbm}
\usepackage{dsfont}
\usepackage{xspace}
\usepackage{cancel}
\usepackage{mathtools}
\usepackage{amsthm}
\makeatletter
\@ifundefined{corollary}{\newtheorem{corollary}{Corollary}}{}
\@ifundefined{lemma}{\newtheorem{lemma}{Lemma}}{}
\makeatother
\newcommand{\E}{\mathbb{E}}
\newcommand{\mc}{\mathcal}
\title[Source-Condition Analysis of Kernel Adversarial Estimators]{Source-Condition Analysis of Kernel Adversarial Estimators}
\author{Antonio Olivas-Martinez, Andrea Rotnitzky}
\date{Source: arXiv:2508.17181v1 (CC BY 4.0)}
\begin{document}
\maketitle

\noindent\emph{Source and licence.} Source-Condition Analysis of Kernel Adversarial Estimators. Version arXiv:2508.17181v1 (CC BY 4.0), distributed under the Creative Commons Attribution 4.0 International licence, as recorded in the arXiv licence metadata for this exact version and shown on the abstract page https://arxiv.org/abs/2508.17181v1. Full rights notice: licenses/arxiv-2508.17181-CC-BY-4.0.txt.\par\smallskip

\noindent\textit{Editorial scope.} Counted unit: the Corollary whose source label is \texttt{corollary:countable\_Lf} in the article (source0.tex lines 1138--1153, source label \texttt{corollary:countable\_Lf}), delivered together with the article's own complete original proof of it (source0.tex lines 1155--1173, one \texttt{proof} environment of 19 lines). No mathematics is written, repaired or re-derived: statement and proof are the article's own text line for line. Required context retained uncounted: lemma:countable (lines 1117--1126). Every definition, display and label of those lines is retained unchanged. Omitted from this package and not redistributed: 1--1116, 1127--1137, 1154--1154, 1174--1440. Nothing inside a retained excerpt is omitted or altered: every retained body is delivered exactly as the article's lines are written, so no omission receipt of any kind accompanies this package. Citations: the retained excerpts carry no citation command at all, so no bibliography entry of the article is transcribed and no reference is redistributed. Reference adaptation: none was needed and none was made; every reference inside the delivered unit resolves inside it, so no unresolved reference or citation is produced. Printed slips kept literal: no printed word, symbol, hypothesis or number of the counted statement or of its proof was altered. Layout and class: the article's own class is \texttt{article} with packages including hyperref, setspace, geometry, amsmath, amsthm, amssymb, scrextend, enumerate, mathrsfs, fancyhdr, tikz, pgfplots, natbib, threeparttable; the wrapper is the lane's pinned \texttt{amsart} editorial wrapper at 10pt on A4, which restates the theorem-like environments the retained text needs and adds the article's own macro definitions that the retained text needs (\texttt{\textbackslash E}, \texttt{\textbackslash mc}), with extra wrapper packages (bm, bbm, dsfont, xspace, cancel, mathtools). The delivered numbering is the unit's own. Assets: no retained span contains a figure environment, a graphics-inclusion command, a PDF-page inclusion, a table or a TikZ picture; no image file is copied into this package, the package declares no asset dependency and no image file is redistributed with it. Licence: version arXiv:2508.17181v1 of this article is distributed under the Creative Commons Attribution 4.0 International licence, as recorded in the arXiv licence metadata for this exact version and shown on the abstract page https://arxiv.org/abs/2508.17181v1. A full rights notice is delivered at licenses/arxiv-2508.17181-CC-BY-4.0.txt. Characters: every retained excerpt is pure ASCII, so the delivered engine, XeLaTeX with its default Latin Modern text font, reports no missing character.
\begin{lemma}
Suppose that:
\begin{enumerate}
    \item $\mc{F}$ is a separable metric space with metric $d$, and let $\mc{S}\subset\mc{F}$ be a countable dense set,
    \item For each $f\in\mc{F}$, $t_f:\mc{O}\mapsto\mathbb{R}$ is a measurable function,  where $\mc{O}$ is the sample space of a random vector $O$,
    \item For each $o\in\mc{O}$, the map $f\mapsto t_f(o)$ is continuous with respect to $d$ at all $f \in \mc{F}$.
\end{enumerate}
Then for all $o_1,\dots,o_n\in\mc{O}$, $\sup_{f\in\mc{F}}\left\{\frac{1}{n}\sum_{i=1}^nt_f(o_i) \right \} = \sup_{f\in\mc{S}}\left\{\frac{1}{n}\sum_{i=1}^nt_f(o_i)\right\}$ and, consequently, for $O_1,\cdots,O_n$ i.i.d. copies of $O$, $\sup_{f\in\mc{F}}\left\{\frac{1}{n}\sum_{i=1}^nt_f(O_i)\right\}$ is a well-defined random variable.
    \label{lemma:countable}
\end{lemma}

\begin{corollary}
Let $\mc{O}_0\subset\mc{O}$ be a measurable set. Suppose that:
\begin{enumerate}
    \item $\mc{F}$ is a separable metric space with metric denoted by $d$ such that $\mc{F}\subset\mc{L}^2(P_O)$, and let $\mc{S}$ denote a countable dense subset of it,
    \item For each $f\in\mc{L}^2(P_O)$, $t_f:\mc{O}\mapsto\mathbb{R}$ is a measurable function,  where $\mc{O}$ is the sample space of a random vector $O$,
    \item For each $o\in\mc{O}$, the map $f\mapsto t_{I_{\mc{O}_0}\cdot f}(o)$ is continuous with respect to $d$ at all $f \in \mc{F}$.
    \item The map $f\mapsto t_{I_{\mc{O}_0}\cdot f}$ is continuous from $(\mc{F},d)$ into $\mc{L}^1(P_O)$.
\end{enumerate}
For given fixed $\beta\in\mathbb{R}$ and any $o\in\mc{O}$, define the map $f\in\mc{F}\mapsto h_f(o):= t_{I_{\mc{O}_0}\cdot f}(o)-\beta\cdot\E\{t_{I_{\mc{O}_0}\cdot f}(O)\}$. Given $\mc{F}'\subset\mc{F}$, define $S':=\mc{F}'\cap\mc{S}$. Then,
\begin{enumerate}
    \item[i)] $\sup_{f\in \mc{F}'}\left\vert \frac{1}{n}\sum_{i=1}^nh_f(O_i) \right\vert
= \sup_{f \in \mc{S}'}\left\vert \frac{1}{n}\sum_{i=1}^nh_f(O_i) \right\vert$. In particular, the left-hand side is a well-defined random variable.
\item[ii)] When $\beta=0$, the conclusion in i) holds even if assumption 4 fails.
\end{enumerate}
 \label{corollary:countable_Lf}
\end{corollary}

\begin{proof} Fix $o\in\mc{O}$. First, we show that the map $f\mapsto h_f(o)$ is continuous with respect to the metric $d$. Take $f_1,f_2\in\mc{F}$. By triangle inequality, we have $ \vert h_{f_1}(o)-h_{f_2}(o)\vert\le \vert t_{I_{\mc{O}_0}\cdot f_1}(o)-t_{I_{\mc{O}_0}\cdot f_2}(o)\vert +|\beta|\cdot\E_P\{\vert  t_{I_{\mc{O}_0}\cdot f_1}(O)-t_{I_{\mc{O}_0}\cdot f_2}(O)\vert\}$.

Let $\varepsilon>0$. By assumption 3, there exists $\delta>0$ such that if $d(f_1,f_2)<\delta$, then $\left\vert t_{I_{\mc{O}_0}\cdot f_1}(o)-t_{I_{\mc{O}_0}\cdot f_2}(o)\right\vert <\varepsilon/2$. By assumption 4, there exists $\delta'>0$ such that  if $d(f_1,f_2)<\delta'$, then $|\beta|\cdot\E_P\{\vert  t_{I_{\mc{O}_0}\cdot f_1}(O)-t_{I_{\mc{O}_0}\cdot f_2}(O)\vert\}<\varepsilon/2$. Then, if $d(f_1,f_2)<\min\{\delta,\delta'\}$, we have $ \vert h_{f_1}(o)-h_{f_2}(o)\vert<\varepsilon$.

By Assumption 1 and definition of $\mc{F}'$, 
$\mc{F}'$ is separable with respect to $d$ with dense subset $\mc{S}'$. Then, it follows from Lemma~\ref{lemma:countable} that 
\[
\sup_{f\in \mc{F}'} \Big\{\frac{1}{n}\sum_{i=1}^nh_f(O_i) \Big\}
= \sup_{f \in \mc{S}'} \Big\{\frac{1}{n}\sum_{i=1}^nh_f(O_i) \Big\},\: \sup_{f\in \mc{F}'} \Big\{-\frac{1}{n}\sum_{i=1}^nh_f(O_i) \Big\}
= \sup_{f \in \mc{S}'} \Big\{-\frac{1}{n}\sum_{i=1}^nh_f(O_i) \Big\},
\]
and the left hand sides of both equalities are well-defined random variables. Therefore
\begin{align*}
    \sup_{f\in \mc{F}'} \Big\vert\frac{1}{n}\sum_{i=1}^nh_f&(O_i) \Big\vert = \max\Bigg\{\sup_{f\in \mc{F}'} \Big[\frac{1}{n}\sum_{i=1}^nh_f(O_i) \Big],\sup_{f\in \mc{F}'} \Big[-\frac{1}{n}\sum_{i=1}^nh_f(O_i) \Big]\Bigg\}
    \\
    &=\max\Bigg\{\sup_{f\in \mc{S}'} \Big[\frac{1}{n}\sum_{i=1}^nh_f(O_i) \Big],\sup_{f\in \mc{S}'} \Big[-\frac{1}{n}\sum_{i=1}^nh_f(O_i) \Big]\Bigg\}=\sup_{f\in \mc{S}'} \Big\vert\frac{1}{n}\sum_{i=1}^nh_f(O_i) \Big\vert
\end{align*}
and the leftmost hand side is a well-defined random variable. This concludes the proof.
\end{proof}

\end{document}
